\chapter{Дополнительные результаты экспериментов}
\label{app:experiments}

В приложении собраны два ablation-эксперимента на смеси CIFAR-100/SVHN
($K=2$, $M=30$ семантических кластеров), на которые ссылается
глава~\ref{ch:experiments}. Все числа --- точность на тесте при кластерном
шлюзе (среднее $\pm$ стд.\ по указанному числу сидов), измеренная по тому
же протоколу, что и таблица~\ref{tab:cifar_main}.

\section{Декомпозиция точности при кластерном шлюзе}
\label{app:decomp}

Чтобы отнести выигрыш MoE к двум его составляющим --- маршрутизации
кластер$\to$эксперт и архитектуре экспертов, --- меняем по одному фактору
за раз (таблица~\ref{tab:decomp}). Блок~(i) фиксирует одну
DARTS-архитектуру (найденную на всей смеси) и меняет только
маршрутизацию: переход от случайного разбиения кластеров к оракульному
разбиению по источнику повышает точность на $3{,}0$ процентных пункта,
то есть маршрутизация --- основной рычаг. Блок~(ii) фиксирует оракульную
маршрутизацию и меняет архитектуру: одна fullmix-архитектура,
CIFAR-найденная и SVHN-найденная различаются в пределах $1$~п.п., то есть
при уже правильной маршрутизации архитектурная гетерогенность экспертов
добавляет мало. Блок~(iii) изолирует вклад \emph{обучения} маршрутизации
внутри SGEM: заморозка матрицы маршрутизации на случайной инициализации
(без M-шага по $\mathbf{R}$) опускает долю совпадения с источником до
случайной ($59\%$), а точность --- до $0.598$, тогда как полный метод
восстанавливает разбиение ($92\%$) и достигает $0.620$ --- по сути того
же значения, что и при оракульной маршрутизации в блоке~(i), но
\emph{без} меток источника.

\begin{table}[h]
\centering
\caption{Ablation, раскладывающий точность при кластерном шлюзе.
«совп.\ с источ.» --- доля совпадения жёсткой маршрутизации с эталонным
разбиением по источнику; она равна $100\%$ по построению всякий раз,
когда используется оракульный сплит.}
\label{tab:decomp}
\begin{tabular}{lcc}
\toprule
Сетап & точность & совп.\ с источ. \\
\midrule
\multicolumn{3}{l}{\textit{(i) меняем маршрутизацию, архитектура фикс.\ (fullmix DARTS)}} \\
\quad случайное разбиение & $0.587 \pm 0.006$ & $\approx 54\%$ \\
\quad оракульное разбиение по источнику & $0.617 \pm 0.007$ & $100\%$ \\
\midrule
\multicolumn{3}{l}{\textit{(ii) меняем архитектуру, оракульная маршрутизация}} \\
\quad fullmix-архитектура & $0.627 \pm 0.009$ & $100\%$ \\
\quad CIFAR-найденная архитектура & $0.636 \pm 0.003$ & $100\%$ \\
\quad SVHN-найденная архитектура & $0.626 \pm 0.006$ & $100\%$ \\
\midrule
\multicolumn{3}{l}{\textit{(iii) обучаем маршрутизацию (SGEM, без меток)}} \\
\quad маршрутизация заморожена (случайно) & $0.598 \pm 0.003$ & $59 \pm 7\%$ \\
\quad \textbf{SGEM (полный)} & $\mathbf{0.620 \pm 0.009}$ & $\mathbf{92 \pm 2\%}$ \\
\bottomrule
\end{tabular}
\end{table}

\section{Чувствительность к весу балансировки нагрузки}
\label{app:lambda}

В главе~\ref{ch:experiments} вес балансировки удерживается малым
($\lambda_{\text{LB}}=1$), потому что правильное разбиение по источнику
само несбалансировано ($19/11$), и сильный штраф балансировки уводил бы
маршрутизацию от него. Сделаем этот компромисс количественным, оценив
три кандидатных жёстких разбиения по \emph{пенализованному} объективу
$L - \lambda_{\text{LB}}\cdot\mathrm{LB}_N$, фактически оптимизируемому на
M-шаге, где пообъектные потери $u_k$ получены реальным обучением
фиксированной архитектуры на кластерах каждого эксперта (без суррогата),
$L = -\sum_k N_k u_k$ --- результирующий взвешенный объектив, а
$\mathrm{LB}_N = K\sum_k P_k^2\cdot N$ --- экстенсивный штраф балансировки
(таблица~\ref{tab:lambda}).

\begin{table}[h]
\centering
\caption{Пенализованный объектив трёх жёстких разбиений на реально
обученных пообъектных потерях. Пенализованная оценка ---
$L - \lambda_{\text{LB}}\,\mathrm{LB}_N$.}
\label{tab:lambda}
\begin{tabular}{lccc}
\toprule
жёсткое разбиение & $u_k$ (по экспертам) & $L$ & $\mathrm{LB}_N$ \\
\midrule
оракул $[19, 11]$ & $[2.73,\ 0.51]$ & $-48099$ & $31491$ \\
баланс, SGEM $[15, 15]$ & $[1.91,\ 1.90]$ & $-55973$ & $29400$ \\
баланс, случайное $[15, 15]$ & $[1.77,\ 2.15]$ & $-57726$ & $29400$ \\
\bottomrule
\end{tabular}
\end{table}

На непенализованном объективе ($\lambda_{\text{LB}}=0$) оракульное
разбиение предпочтительнее с запасом $\Delta L = 7874$. Однако, будучи
несбалансированным, оно штрафуется сильнее балансированного на
$\Delta\,\mathrm{LB}_N = 2091$ на единицу $\lambda_{\text{LB}}$. Поэтому
оракульное разбиение выигрывает пенализованный объектив лишь при
$\lambda_{\text{LB}} < \Delta L / \Delta\,\mathrm{LB}_N \approx 3.77$; за
этой точкой перелома балансированное (и потому доменно-слепое) разбиение
получает более высокую оценку. Это прямой, не зависящий от суррогата
конфликт между балансировкой нагрузки и несбалансированной доменной
структурой; он объясняет, почему большое $\lambda_{\text{LB}}$ ---
надёжно предотвращающее схлопывание --- подавляло бы и саму искомую
специализацию, поэтому мы работаем существенно ниже точки перелома.
